\section{Equivariant Hecke systems and specialization}\label{sec:conventions}

We fix the sheaf conventions and explain which parts of the geometric
specialization formalism carry a semilinear action. The issue is compatibility
of the entire Hecke system, including collisions of moving points.

\subsection{Sheaves, level structures, and moving legs}

Write $E=\E$, $G=\SL_n$, and $\Gd=\PGL_n$, with the split forms understood
over every base field. On a smooth stack $\mathcal B$ locally of finite type,
$D_{\lcc}(\mathcal B,E)$ denotes the category of geometric adic complexes
whose pullback to each smooth finite-type scheme chart is bounded and
constructible. The bounds may depend on the chart. A complex $F$ is perverse
if $b^*F[d]$ is perverse for every smooth chart $b:Y\to\mathcal B$ of
constant relative dimension $d$. These conventions agree with
\cite[Section~2.1]{CT}.

For a pointed smooth projective curve $(X,x)$, put
\[
 \mathcal A=\Bun_{G,B,x}(X),\qquad
 \mathcal A^+=\Bun_{G,R_u(B),x}(X),\qquad T=B/R_u(B).
\]
An enhanced flag is a reduction to $R_u(B)$; forgetting the enhancement is
a representable $T$-torsor $\mathcal A^+\to\mathcal A$.
At ordinary Borel level, the trace form identifies a cotangent vector with
a Higgs field
\[
 \theta\in H^0\bigl(X,\ad(P)\otimes\omega_X(x)\bigr),
 \qquad \operatorname{Res}_x\theta\in\operatorname{Lie}R_u(B_P),
\]
where $B_P$ is the chosen Borel in the fiber at $x$. The cone
$\Lambda_{\mathrm{par}}$ consists of those fields that are nilpotent over
the function field of $X$. On a smooth chart, singular support means the
geometric adic singular support of the pulled-back complex, contained in
the smooth cotangent pullback of this cone; see
\cite[Section~3.1]{CT} and \cite{Beilinson,Barrett}.

Let $U$ be the locus of allowed modifications, disjoint from the level
point. For a finite set $I$, the Hecke stack $\mathcal H_I$ has an input
bundle map $p_I$ and an output-and-leg map
$o_I:\mathcal H_I\to\mathcal B\times U^I$.
If $\boldsymbol V=(V_i)_{i\in I}$ are representations of $\Gd$, define
\begin{equation}\label{eq:hecke-normalization}
 \Hecke_{I,\boldsymbol V}(F)
  =o_{I,*}\bigl(F\,\widetilde\boxtimes\,
                      \operatorname{Sat}_I(\boldsymbol V)\bigr).
\end{equation}
The twisted product is ordinary exterior product in input frames, descended
using equivariance of the Satake kernel. All pushforwards are taken on
finite Schubert bounds, where the output maps are proper. On a fixed-point
Schubert orbit of dimension $d_\lambda=\langle2\rho_G,\lambda\rangle$,
the simple kernel is normalized by
\begin{equation}\label{eq:satake-normalization}
 \IC_\lambda|_{\Gr^\lambda}=E[d_\lambda](d_\lambda/2).
\end{equation}
Here $\rho_G$ is the half-sum of positive roots. Since the cocharacter
lattice of $G$ is its coroot lattice, $d_\lambda$ is even. In particular
all twists in \eqref{eq:satake-normalization} are integral.
There is no moving-leg shift in \eqref{eq:hecke-normalization}.
We use geometric Satake with its fusion tensor structure \cite{MV}.

For a surjection $a:I\twoheadrightarrow J$, let
$\delta_a:U^J\to U^I$ be the collision diagonal and set
$W_j=\bigotimes_{i\in a^{-1}(j)}V_i$. The relative fusion convention is
\begin{equation}\label{eq:fusion-convention}
 (1\times\delta_a)^*\Hecke_{I,\boldsymbol V}(F)
   \simeq\Hecke_{J,\boldsymbol W}(F).
\end{equation}
This is ordinary pullback, without a codimension shift.

\begin{definition}\label{def:eigenstructure}
For a $\Gd$-local system $\vartheta$ on $U$, a coherent Hecke
eigenstructure on $F$ consists of isomorphisms
\begin{equation}\label{eq:eigenstructure}
 \Hecke_{I,\boldsymbol V}(F)
   \simeq F\boxtimes\bigboxtimes_{i\in I}(V_i)_\vartheta
\end{equation}
for all finite $I$, natural in the representations, and compatible with
the unit, successive modifications and convolution, permutations, and
all collision maps \eqref{eq:fusion-convention}, including their iterated
compatibilities. The notation $(V)_\vartheta$ means the lisse sheaf
obtained by evaluating the representation $V$ on $\vartheta$.
\end{definition}

Let $\sigma$ be a semilinear automorphism of the ground field and the curve,
preserving the split group and the level data. An equivariant sheaf has an
isomorphism $\sigma^*F\simeq F$; an equivariant eigenstructure requires
\eqref{eq:eigenstructure} to commute with this isomorphism and the specified
isomorphism $\sigma^*\vartheta\simeq\vartheta$. We use the group generated
by $\sigma$, so the one isomorphism determines all its iterates and inverses.
Transport acts on the base, not on the coefficient field $E$.
At a finite-field fiber the ground-field action $a\mapsto a^q$ gives the
arithmetic convention. Inverting the generator everywhere gives the usual
geometric Frobenius convention for Weil sheaves.

\subsection{Satake transport and its Weil normalization}

\begin{lemma}\label{lem:satake-transport}
Semilinear isomorphisms of pointed curves identify the relative Satake and
Hecke systems, compatibly with all the operations in
Definition~\ref{def:eigenstructure}. In representation labels the required
symmetric tensor comparison is unique. The same assertion holds for
extension between algebraically closed characteristic-zero fields.
\end{lemma}

\begin{proof}
Transport preserves the split Schubert orbits, their intersection complexes,
convolution diagrams, and fusion diagrams. Thus it gives a symmetric tensor
equivalence of the spherical Satake categories preserving all highest-weight
labels. Under Satake it is a symmetric tensor autoequivalence of
$\Rep_E(\PGL_n)$ with trivial outer automorphism. Over the algebraically
closed coefficient field it is tensor isomorphic to the identity: the
fiber-functor torsor is trivial and the remaining group automorphism is inner.
Two such isomorphisms differ by a tensor automorphism of the identity functor,
namely an element of $Z(\PGL_n)$, which is trivial.

Use this comparison for the full geometric system. On separated legs it is
the exterior product comparison. The fusion extensions, proper convolution
maps, and their exchange transformations are transported together, so their
compatibility maps agree. In the kernel categories the fusion extensions
are middle extensions with the usual perverse leg shifts; a comparison on
the separated locus determines them. Applying the corresponding natural
sheaf operations gives the Hecke comparison. Algebraically closed base
extension has the same properties: it preserves the simple kernels and
convolution and is an equivalence on spherical Satake. Uniqueness again
identifies its representation labels.
\end{proof}

\begin{lemma}\label{lem:normalized-weil-satake}
Over a finite field, the split Weil structures on
\eqref{eq:satake-normalization} give a symmetric tensor Satake system,
with ordinary coefficient multiplicity spaces carrying trivial action.
Its Frobenius comparison is the one in
Lemma~\ref{lem:satake-transport}.
\end{lemma}

\begin{proof}
We verify the normalization, since purity alone would not specify the
Frobenius scalars. Resolve a split affine Schubert variety by its
Bott--Samelson resolution, using the Iwahori stratification. Its iterated
projective-line-bundle structure is defined over the finite field. The
projective-bundle formula gives scalar geometric Frobenius $q^j$ on
$H^{2j}$ and zero odd cohomology. After the shift and twist
$[d_\lambda](d_\lambda/2)$, the scalar on degree $i$ is $q^{i/2}$.

In perverse degree zero of the proper direct image, the full-support
intersection complex occurs with multiplicity one. The decomposition
theorem \cite[Theorems~5.3.8 and~5.4.5]{BBD} makes its isotypic summand geometrically well defined and hence
Frobenius stable, with the normalization fixed on the open orbit. The
Frobenius-stable perverse Leray filtration, whose spectral sequence
degenerates by the same theorem, exhibits each cohomology group of this
summand as a subquotient of the corresponding scalar cohomology group.
Thus
\[
 H^i(\Gr,\IC_\lambda)\text{ has Frobenius as on }E(-i/2)
 \quad(i\text{ even}),\qquad H^i=0\quad(i\text{ odd}).
\]

The cohomological convolution comparison is graded and Frobenius
compatible. Indeed it is constructed by K\"unneth on separated legs and
fusion over the split affine line
\cite[Lemma~6.1, Equations~(6.2a)--(6.2c), and Section~14]{MV}.
The cohomology system is lisse across
collisions and geometrically constant: it is constant on the separated
locus in the coordinate trivializations, and lisseness extends this
constancy. Comparison between rational configurations therefore commutes
with Frobenius; binary tensor comparison uses the rational configurations
$(0,1)$ and $(0,0)$, available over every finite field.
The scalar rule is consequently the same for convolution
and the corresponding tensor product. Every graded tensor comparison on
total cohomology commutes with these scalars. Faithfulness of the
cohomological Satake fiber functor shows that the kernel tensor maps do
as well. In particular, the induced Frobenius action on the geometric
multiplicity space $\Hom(\IC_\nu,\IC_\lambda*\IC_\mu)$ is trivial:
source and target cohomology have the same scalar in each degree,
and the cohomological fiber functor is faithful.
Fusion extends the compatibility to the relative system.
Finally, the uniqueness in Lemma~\ref{lem:satake-transport} identifies
this Weil comparison with the transport comparison.
\end{proof}

\subsection{The specialization input with an action}

The geometric result we need is stronger than one-point Hecke commutation.
It is the combination of Propositions~2.1 and~5.1 of \cite{CT}, with
the following relative setup. Let $S=\Spec R$ be an excellent complete
strictly henselian discrete valuation trait with algebraically closed
residue field and with $\ell$ invertible. Fix a geometric generic point
$\bar\eta$ and write $s$ for the closed point. Let $\mathcal B/S$ be a
smooth locally finite-type bundle stack and let $U/S$ be a smooth scheme
of relative dimension one of allowed legs. Bounded Hecke correspondences
have proper representable output maps, and in coordinates and input frames
they have the split Schubert local models with their spherical kernels.
The setup includes ordinary and enhanced level disjoint from the legs,
and the twisted nodal families used below.

\begin{proposition}[Geometric specialization and transport]
\label{prop:equivariant-specialization}
In this setup, geometric nearby cycles preserve local bounded
constructibility and are perverse exact. They admit natural isomorphisms
\begin{equation}\label{eq:hecke-specialization}
 \Hecke^s_{I,\boldsymbol V}(\Psi F)
  \xrightarrow{\ \sim\ }
  \Psi\bigl(\Hecke^{\bar\eta}_{I,\boldsymbol V}(F)\bigr)
\end{equation}
compatible with the unit, convolution, permutations, and every fusion
diagonal. Suppose $F$ has a coherent eigenstructure with parameter
$\vartheta_{\bar\eta}$. For each representation, suppose its evaluation
sheaf has a lisse lattice whose finite reductions extend lisse after
finite trait extensions, compatibly on further extensions and with the
reductions of a special-fiber tensor system $\vartheta_s$. Then $\Psi F$
has a coherent eigenstructure with eigenvalue $\vartheta_s$.

If the models, $F$ with its given coherent eigenstructure, and these
tensor-specialization data carry a semilinear generator, with a
compatible lift to $\bar\eta$, all the
comparisons and the resulting eigenstructure are equivariant. For the
finite-stage lattice models one may use compatible invariant trait
extensions. No single finite extension is required for all reductions.
\end{proposition}

\begin{proof}
The geometric assertions are exactly
\cite[Propositions~2.1 and~5.1]{CT}; the foundational nearby-cycle
formalism is also described in \cite[Section~4.1]{GR} and
\cite[Theorems~4.1 and~6.7]{HansenScholze}. In particular, these are
geometric nearby cycles without inertia invariants, and the input need
not descend over one finite extension of the fraction field.

We explain the additional equivariance. An isomorphism of traits and
models, together with its lift to the geometric generic point, identifies
the defining nearby-cycle diagrams. The pullback, proper-pushforward,
and tensor exchange maps therefore commute with its transport. In input
frames the Hecke comparison is the natural exterior-product exchange
with a Satake kernel, followed by proper-pushforward exchange. This is
the construction in \cite[Section~5]{CT}, including its comparison on
collision diagonals. The Satake-label identification in
\cite[Lemma~5.3]{CT} also commutes with transport: the two composites
are tensor comparisons and hence agree by
Lemma~\ref{lem:satake-transport}.

On the eigenvalue side, the lisse tensor comparison is natural at every
finite lattice reduction and every further trait extension. Passing to
the adic system preserves this naturality. Thus
\eqref{eq:hecke-specialization} transports the equivariant eigenmaps
and all their diagrams. The convolution and nested diagonal identities
are identities of the same exchange transformations; they require no
separate choices. At a finite-field fiber their Satake structures are
the normalized Weil structures by
Lemma~\ref{lem:normalized-weil-satake}.
\end{proof}

\subsection{The characteristic restriction for \texorpdfstring{$\SL_n$}{SLn}}

The remaining geometric inputs of \cite{CT} have four Lie-theoretic
hypotheses: a nondegenerate invariant symmetric form, also nondegenerate
on every Levi's Lie-algebra center; Chevalley restriction for every Levi;
scheme-theoretic semisimple centralizers that are Levis; and containment
of every nilpotent element, over every field extension, in the nilradical
of a parabolic defined over that field. We verify them to make the
characteristic restriction in the main theorem explicit.

\begin{lemma}\label{lem:lie-hypotheses}
For $G=\SL_n$ over an algebraically closed field of characteristic zero or of characteristic
$p>n$, the hypotheses H1--H4 of \cite[Theorem~1.1]{CT} hold.
\end{lemma}

\begin{proof}
The invariant form $\kappa(X,Y)=\operatorname{tr}(XY)$ is nondegenerate
on $\mathfrak{sl}_n$, since $n$ is invertible. For a block Levi with
block sizes $n_1,\ldots,n_r$, its restriction to the Lie-algebra center
is the diagonal form with entries $n_i$ on
\[
 \{(z_i):\textstyle\sum_i n_i z_i=0\}.
\]
This hyperplane is the orthogonal complement of $(1,\ldots,1)$,
whose squared length is $n$. All $n_i$ and $n$ are invertible, so the
restricted form is nondegenerate.

Chevalley restriction for these Levis follows by block diagonalization
on the dense regular-semisimple locus. Block-symmetric polynomials in
the eigenvalues lift by the block characteristic polynomials.
The Weyl-group order divides $n!$, which is invertible; averaging shows
that invariants on the total-trace-zero hyperplane are restrictions of
invariants on the full diagonal space. Block diagonalization may be
performed in the determinant-one Levi, by adjusting the determinant
with an element of the diagonal centralizer.

The scheme-theoretic centralizer of a semisimple Lie-algebra element
is its block Levi, as follows directly from the matrix commutation
equations. Finally, over every field extension, vanishing of the
positive-degree invariant polynomials says that the matrix is nilpotent.
A basis adapted to its kernel filtration puts it in a strictly triangular
Lie algebra defined over that field. It therefore lies in the
nilradical of a rational Borel, which supplies the required rational
parabolic.
\end{proof}
